\answer{ex:07:1}{$V=\kappa p/\bigl(\kappa p+(1-\kappa)(1-p)\bigr)$: $0.059$، $0.20$، $0.50$؛ الوسيط يساوي $0$، أما النقطة~\eqref{eq:expectile} فتتغير مع $\kappa$ تغيرًا سلسًا.}
\answer{ex:07:2}{$p=1/3$، $x=0.75$، $\sigma_r=0.35$، $V(0.2)=0.083$.}
\answer{ex:07:3}{$V=\sqrt\kappa/(\sqrt\kappa+\sqrt{1-\kappa})$، من $0.25$ إلى $0.75$؛ التشتت $\propto\sqrt\eta$؛ وللقطرتين $V=0.25+0.5\kappa$، فيختلف التوزيعان بـ$0.05$ عند العصبونات الطرفية، ويلزم $\eta\lesssim0.01$.}
\answer{ex:07:4}{(أ)~$J^*=2\sqrt{C\bar R}$. (ب)~الزيادة $\bigl(k^{1/2}+k^{-1/2}\bigr)/2-1$: $6\,\%$ عند $k=2$ و$25\,\%$ عند $k=4$، فتكفي دقة ”في حدود الضعف“.}
\answer{ex:07:5}{ذروة مساحتها $R(e^{-1}-e^{-2})\approx0.23R$، أي ما يعادل $2.3\,\text{s}$ من الارتفاع التدريجي؛ وإذا كان الدوبامين يخبر بـ$V$ فلا ذروة، بل يقفز المستوى درجةً فيكبر $e$ مرة.}
\answer{ex:07:6}{الحدث يزيح الوزن $\propto A\tau_f/(\tau_e+\tau_f)$، والمستوى يعطي خطأ $s\tau_f$: $9\,\text{s}\le\tau_f\le17\,\text{s}$.}
\answer{ex:07:7}{في البداية $\Delta P=0.80$، $M=0.90$. (أ)~$\Delta P=0.74$ ($-8\,\%$)، $M=0.43$ ($-52\,\%$)؛ (ب)~$\Delta P=0.40$ ($-50\,\%$)، $M=0.80$ ($-11\,\%$): تفارق مزدوج.}
\answer{ex:07:8}{$N_{\text{eff}}\to2+\rho^2N\approx10^6$ محاولة: أقل بـ$10^4$ مرة مما مع سلك واحد، لكن تسربًا قدره $1\,\%$ ما زال يكلّف مليون محاولة.}
